\documentclass{article}
\usepackage{../Self_Style}

\usepackage{tikz-cd} %as a current substitution of xypic...

\usepackage[normalem]{ulem} %for strike through

\title{Equivariance in Category Theory}
\author{Zih-Yu Hsieh}
\date{\today}

\begin{document}
\maketitle
\tableofcontents

\section{Motivation}

To study some ``things'', what's a good way of studying it? One thing to do, is studying some other things that're interconnected to it.

For instance, in topological K-theory, people study vector bundles over a space $X$, and turn them into sequence of spaces called ``spectrum'' associated to $X$. This sequence of spaces and their topologies gives some interesting invariants of the space itself.

Similarly, in algebraic K-theory, people study ``projective modules'' over a ring $R$, we can also construct spectrum over this ring, and the topologies also give some interesting invariants of the ring.

\hfil

However, for both of these cases (vector bundles and projective modules), they're both closed under direct sums and tensor products (whatever that is), so we actually starts with a ``category'' with some nice product structure. These are called symmetric monoidal categories. In general, a K-theory can be done, by turning some symmetric monoidal category into a spectrum. So, classifying them turns out to be extremely important.

Now, analogous constructions can be done if the ``things'' endow some kind of group actions, which is the main study of equivariant K-theory. For this, an important question must be asked: What the heck is a ``Symmetric Monoidal $G$-Category''? Today, we try to get to this, by first introducing equivariance in categories.

\hfil

\section{What is Symmetric Monoidal Category?}

Now we'll switch gear, and talk about three non-equivariant categories:

\begin{defn}
    A category $\A$ is a \emph{symmetric moniodal caetgory}, if there is a binary product $\tensor:\A\times\A\rightarrow\A$, that is unital, associative, and commutative up to natural isomorphism.
\end{defn}

Which, we collect them (and functors preserving the product structure) as $\SymMon$.

Next, we're gonna blackbox whatever an operad is. Yes. (Also, there are many texts by Peter describing this). Anyway, a ``slogan'' for operad is, that each $n\in\NN$ has some ``things'' that encodes how to define $n$-fold products on some ``other things'', and how to combine these products.

In the category of small categories $\Cat$ (with functors as 1-morphisms, natural transformation as 2-morphisms), we have an operad $\cP$, such that a notion of ``pseudoalgebra'' over this $\cP$, precisely describes a category $\A$ and a collection of $n$-fold products on the category (using a sequence of functors $\cP(n)\times \A^n\rightarrow\A$). Collect tese pseudoalgebras as $\cP$-$\PsAlg$.

Also, if collecting the category of based finite sets $\F$, there are also some ``pseudofunctors'' $\F\rightarrow\Cat$ (whatever that is) which describes a category $\A$ and a collection of $n$-fold products. Collect these pseudofunctors as $\F_\RR$-$\PsAlg$. 

These are some imputs to various stages of an infinite loop space machine (whatever that is). So an interesting question is: are they the same? I fact, they are!

\begin{thm}
    The 2-categories $\SymMon$, $\cP$-$\PsAlg$, and $\F_\RR$-$\PsAlg$ are equivalent.
\end{thm}

\begin{proof}(Really Rough Sketch, skip if needed)
    
    Given a symmetric monoidal category $\V$, for each $n$, fix a representative of $n$-fold product (as they're all isommorphic), it constructs a functor $\cP(n)\times\V^n\rightarrow\V$, and turns out to satisfy the operadic pseudoalgebra's property (modulo a bunch of details).

    Given a $\cP$-pseudoalgebra $\A$, there is $\cP(2)\times \A^2\rightarrow\A$, which helps us define a binary product. The requirement on permutation and ways of combining it to $n$-fold product $(n\geq 3)$ describes commutativity and associativity (modulo a bunch of details again). Plus, we have $\cP(0)\rightarrow \A$ (with $\cP(0)=*$), this specifies our unit.

    \hfil

    Similarly, given $\A$ a $\cP$-pseudoalgebra, for every based finite set $\mathbf{n}=\{0,1,...,n\}$, associate it to $\A^n$, then every based function $f:\mathbf{n}\rightarrow\mathbf{m}$ can be translated to $\A^n\rightarrow\A^m$ by mapping $\A^{|f^{-1}(i)|}\rightarrow\A$ on the $i$th coordinates, and empty coordinates specified using the unit $*\rightarrow \A$ (given by the pseudoalgebra structure).

    Conversely, any $\F_\RR$-pseudoalgebra consists of the $n$th space as $\A^n$ (where $\A$ is the $1$st space). Then, the canonical $n$-fold product $\phi_n:\mathbf{n}\rightarrow\mathbf{1}$ ($i\mapsto 1,0\mapsto 0$) defines an $n$-fold product $\A^n\rightarrow\A$, and pseudofunctoriality comes into play to define $\cP$-pseudoalgebra structure.
\end{proof}

\hfil

\section{$G$-Categories and 2-Category $\Cat_G$}

Fix a group $G$ (say finite).

\begin{defn}
A small category $\A$ is a (strict) \emph{$G$-category}, if there exists a map $H:G\rightarrow \Cat(\A,\A)$, such that $H(e)=\id_\A$, and $H(gh)=H(g)\circ H(h)$.
\end{defn}

So, each group element is a functor, and an automorphism on $\A$, where composition functors are compatible with group structures. Just like how it's defined for $G$-sets.

\begin{defn}
    The 2-category $\Cat_G$ consists of all small $G$-categories, the 1-morphisms between $G$-categories are functors (not necessarily equivariant), and the 2-morphisms are natural transformations between them.
\end{defn}

An interesting thing to observe: each $G$-category $\A$ realizes $g\in G$ as an automorphism on $\A$. So, if $F:\A\rightarrow\B$ is a functor, so is $g\circ F\circ g^{-1}$! This realizes the ``functor category'' $\Cat_G(\A,\B)$ as a small $G$-category! 

If consider the subcategory fixed by all functors of $G$, $\Cat_G(\A,\B)^G$ gives all the $G$-equivariant functors from $\A$ to $\B$.

\hfil

\section{What's Symmetric Monoidal $G$-Category?}

\hfil

Now, in $\Cat_G$, there is an operad called $\cP_G$, which is an analous to $\cP$ in $\Cat$. 

\begin{defn}
A \emph{symmetric monoidal $G$-category}, is defined to be $\cP_G$-pseudoalgebra.
\end{defn}

A problem is: we need to describe all the $n$-fold product and equivariant structure, before even studying symmetric monoidal $G$-category, and it's genuinely hard to understand the structure. So what the heck is a symmetric monoidal $G$-category? No one knows in general. If one can describe it using just binary product, it may be a lot easier.

\hfil

As an analogy, we also have the category of based finite $G$-sets, $\F_G$, and also a way of expressing these $\cP_G$-pseudoalgebras as pseudofunctors $\F_G\rightarrow \Cat_G$ (denoted $\F_{G,\RR}$-$\PsAlg$). These are again two categories showing up in equivariant $K$-theory (when constructing $G$-spectrum). After having all the correspondance in non-equivariant case, one can also ask: Are these two damn categories equivalent? If so, HOW?

This is what my project is doing currently, and hope this can also answer the previous question in some way, seeing a better structure in symmetric monoidal $G$-categories.

\end{document}